% TOP_PROBLEM: {"id": 3074, "title": "Big Line or Big Clique in Planar Point Sets", "category_id": 6, "category": {"id": 6, "name": "geometry", "display_name": "Geometry", "description": "Euclidean and non-Euclidean geometry, geometric structures.", "slug": "geometry", "order_index": 6, "created_at": "2026-07-31T15:26:25.671Z"}, "status": "open", "problem_number": "OPG-588", "set_id": 12}
% FIRST_POSTED: 2026-10-01
% ATTEMPT_STATUS: unresolved
% UnsolvedMath ID 3074; source catalogue code OPG-588
% !TeX program = lualatex
% GPT-6 Astra Ultra research attempt; independent partial work, 2026-10-01.
\documentclass[11pt,a4paper]{article}
\usepackage[margin=25mm]{geometry}
\usepackage{fontspec}
\setmainfont{lmroman10-regular.otf}[BoldFont=lmroman10-bold.otf,
 ItalicFont=lmroman10-italic.otf,BoldItalicFont=lmroman10-bolditalic.otf]
\usepackage{amsmath,amssymb,mathtools}
\usepackage{unicode-math}
\setmathfont{latinmodern-math.otf}
\AtBeginDocument{\let\setminus\smallsetminus}
\usepackage{xurl,hyperref}
\hypersetup{colorlinks=true,urlcolor=blue}
\setcounter{secnumdepth}{0}
\setlength{\parindent}{0pt}
\setlength{\parskip}{5pt}
\setlength{\emergencystretch}{3em}
\begin{document}
\section*{Big Line or Big Clique in Planar Point Sets}\label{3074}
{\small UnsolvedMath ID 3074 · OPG-588 · Geometry · OpenGarden}
\subsection{Definitions and mathematical statement}
Let $S\subset\mathbb R^2$ be a finite set of distinct points. For
$x,y\in S$ with $x\ne y$, their open segment is
$(x,y)=\{(1-t)x+ty:0<t<1\}$. Say that $x$ and $y$ are visible with
respect to $S$ if $(x,y)\cap S=\varnothing$. The visibility graph has
vertex set $S$ and joins precisely the visible pairs.

The conjecture is that for all integers $k,\ell\ge2$, there exists an
integer $N(k,\ell)$ such that every finite $S$ with $|S|\ge N(k,\ell)$
satisfies at least one of the following alternatives:
\[
 \exists\text{ line }L:\ |S\cap L|\ge\ell,
 \qquad\text{or}\qquad
 \exists T\subseteq S:\ |T|=k\text{ and every pair in }T\text{ is visible}.
\]
Thus the second alternative is a clique of order $k$ in the visibility
graph. The threshold may depend on the two requested sizes, but not on
the coordinates of $S$.

Visibility is tested against the entire original point set, including
points outside $T$. Merely choosing $k$ points with no three collinear
does not ensure visibility: other points of $S$ may obstruct their
segments. Conversely, pairwise visibility does not require the convex hull
of $T$ to contain no additional points away from those segments.
Collinear configurations are allowed, and there is no general-position
assumption. The assertion is specifically about finite sets of sufficiently
large cardinality; replacing them by an infinite set is a different statement.
\subsection{Short English statement}
Is every sufficiently large finite planar point set forced to contain many collinear points or a large set of pairwise visible points?
\subsection{Research attempt: minimal triangles and the obstruction from blockers}
\paragraph{Scope and literature check.}
GPT-6 Astra Ultra, 1 October 2026. Visibility is always measured
against the full point set. The historical Garden discussion [S2]
cannot be used as a current list of solved cases: Bonnet's August
2026 preprint [S3] proves the four-collinear-or-six-visible case.
The September 2026 paper [S4] gives further near-clique results.
Neither is attributed to this attempt. Here we prove an elementary
small-clique case and audit the standard general-position extraction.

\paragraph{1. A minimum-area triangle is visible.}
Suppose $S$ is finite and not collinear. Among its nondegenerate
triangles choose one $abc$ with minimum positive area, possible because
there are finitely many triples. If $z\in S$ lies strictly between
$a$ and $b$, then $azc$ is nondegenerate and has area
\[
 [azc]=\frac{|az|}{|ab|}[abc]\in(0,[abc]),
\]
contradicting minimality. The same reasoning applies to the other two
sides. Thus $a,b,c$ are pairwise visible in the entire set $S$.
It follows that $N(3,\ell)=\max(3,\ell)$ is a sufficient threshold:
a collinear set of that size has an $\ell$-point line, and every
noncollinear set has the visible triangle just proved.
This is a sufficient threshold, not a claim of its optimality for
every convention at the small endpoints.

\paragraph{2. A large general-position subset can be extracted.}
Assume $\ell\ge3$ and no line contains $\ell$ points of $S$.
Take an inclusion-maximal subset $T\subseteq S$ with no three
collinear, and write $m=|T|$, $n=|S|$. Every point of $S\setminus T$
lies on a line through two points of $T$, or it could be added.
Such a line contains exactly two points of $T$ and at most
$\ell-3$ further points of $S$. Counting over the $\binom m2$
lines, allowing overcounting outside points, gives
\[
 n\le m+(\ell-3)\binom m2.
\]
Consequently a threshold exceeding
$(k-1)+(\ell-3)\binom{k-1}{2}$ guarantees $k$ points with no three
collinear. When $\ell=3$ the entire set is already in general
position, and every pair is visible, so that special case of the
original conjecture is immediate.

\paragraph{3. Why this extraction does not give visibility.}
The preceding argument tests collinearity within $T$, while visibility
tests for blockers in $S\setminus T$ too. For an explicit obstruction,
let $S=\{0,1,2\}^2$ and let $T$ be its four corner points. Each
side of the corner square has a midpoint in $S$ and each diagonal
has the centre in $S$. Thus no pair of points of $T$ is visible,
although no three points of $T$ are collinear. This also blocks the
naive application of graph Ramsey theory: an independent set in the
visibility graph need not be a collinear set. The four corners are
an independent set with several directions.

A convex-position extraction has the same issue for its diagonals
and even its sides. One would need control of the blockers in the
ambient set, rather than only of the geometry of the selected points.
Likewise an empty interior alone does not ensure clear boundary
segments if the definition of emptiness ignores boundary points.

\paragraph{Verification and remaining gap.}
The minimum-area argument uses a strict interior blocker, so the
area decreases strictly; it never excludes unrelated points inside
the triangle. The grid example was checked using the explicit six
pair segments, four side midpoints and the common diagonal midpoint.
It does not disprove the conjecture, since other visible triangles
exist in the same set. The remaining general step is a large
ambiently visible clique under a bounded-collinearity hypothesis;
neither the pair-line count nor generic graph Ramsey theory supplies it.

\paragraph{Reproduction of finite checks.}
The finite checks stated above were run with Python 3 using the repository
script [C1]. They verify the stated examples and finite ranges only.

\subsection{Final status}
UNRESOLVED for arbitrary $k,\ell$. The visible-triangle case and a
general-position bound are proved, and the proposed extraction-to-clique
inference is disproved by an explicit example. Existing 2026 advances
are cited; no claim of new general progress or full completion is made.

\subsection{Sources}
\begin{itemize}
\item[S1] ULAM.AI and the UnsolvedMath Contributors, supplied problems.json, record 3074 (OPG-588), OpenGarden collection. Catalogue identity and source transcription; CC BY 4.0.
\url{https://huggingface.co/datasets/ulamai/UnsolvedMath}.
\item[S2] Open Problem Garden, Big Line or Big Clique in Planar Point Sets. Original problem and discussion; the mathematical formulation below is expanded from the supplied source record.
\url{http://www.openproblemgarden.org/op/big_line_or_big_clique_in_planar_point_sets}.
\item[S3] É. Bonnet, \emph{Large Finite Point Sets Have 4 Collinear Points or a 6-Clique}, arXiv:2608.19468, August 2026. Checked 1 October 2026.
\url{https://arxiv.org/abs/2608.19468}.
\item[S4] B. B. Bhattacharya, S. Das, S. S. Islam, A. Mohapatra and S. Sen, \emph{Large Planar Point Sets Contain 4 Collinear Points or Almost 7-Cliques, and Related Results}, arXiv:2609.25727, September 2026. Checked 1 October 2026.
\url{https://arxiv.org/abs/2609.25727}.
\item[C1] Reproducible finite-check code, executed 1 October 2026 with Python 3 (standard library only):
\path{scripts/check_attempt_batch_geometry_2026_10_01.py}.
\end{itemize}
\end{document}
